\chapter{\ploren{Wyniki i dyskusja}{Results and discussion}}

% Zastąp poniższe wskazówki treścią właściwą dla pracy. Układ można podzielić
% na dwa rozdziały: „Wyniki” i „Dyskusja”, jeżeli materiał jest rozbudowany.

\begin{quote}
\itshape
\ploren
  {Przedstawiony układ jest propozycją. Wyniki i dyskusję można prowadzić łącznie — omawiając każdy rezultat bezpośrednio po jego przedstawieniu — albo rozdzielić na osobne części. W obszernej pracy dopuszczalne są dwa odrębne rozdziały. Niezależnie od układu trzeba wyraźnie odróżnić dane uzyskane w badaniach od ich interpretacji. Tekst tej wskazówki należy usunąć z ostatecznej wersji pracy.}
  {The proposed structure is a recommendation. Results and discussion may be combined, with each result interpreted immediately after presentation, or separated into distinct parts. Two separate chapters are acceptable in an extensive thesis. Regardless of the structure, data obtained from the study must be clearly distinguished from their interpretation. This guidance should be removed from the final version of the thesis.}
\end{quote}

Rozdział powinien odpowiadać na pytanie, co wykazała realizacja projektu lub przeprowadzone badania oraz jakie znaczenie mają uzyskane rezultaty. Kolejność prezentacji najlepiej powiązać z celami szczegółowymi, pytaniami badawczymi, hipotezami albo wymaganiami określonymi wcześniej. Nie należy przedstawiać wyników w kolejności ich przypadkowego uzyskiwania.

\section{\ploren{Sposób opracowania wyników}{Result processing}}

Na początku należy krótko przypomnieć miary, metody agregacji i zasady analizy danych, jeżeli są potrzebne do zrozumienia prezentowanych rezultatów. Szczegółowa procedura powinna znajdować się w poprzednim rozdziale; tutaj wystarczy wskazać sposób wyznaczania wartości przedstawionych na wykresach i w tabelach, na przykład średniej, mediany, odchylenia standardowego, przedziału ufności lub niepewności pomiaru.

Trzeba podać liczbę uwzględnionych obserwacji lub powtórzeń oraz wyjaśnić ewentualne odrzucenie danych. Nie należy dobierać sposobu analizy dopiero po obejrzeniu wyników bez ujawnienia tej decyzji. Jeśli zastosowano dodatkowe przetwarzanie względem wcześniej opisanej metodyki, zmianę należy wyraźnie zaznaczyć.

\section{\ploren{Prezentacja wyników}{Presentation of results}}

Wyniki należy przedstawiać w formie najlepiej odpowiadającej ich charakterowi. Tabela służy przede wszystkim do podawania dokładnych wartości i porównywania wielu parametrów, natomiast wykres powinien uwidaczniać zależności, trendy, rozrzut lub różnice między wariantami. Tych samych danych nie należy bez potrzeby powtarzać jednocześnie w tabeli, na wykresie i w tekście.

Każdy rysunek i każda tabela muszą zostać przywołane oraz omówione w tekście. Opis powinien wskazywać najważniejszy rezultat, a nie przepisywać wszystkie wartości. Osie wykresów wymagają nazw i jednostek, legenda powinna być jednoznaczna, a użyte kolory oraz oznaczenia muszą pozostać czytelne także w druku w skali szarości. Wykresy nie powinny sugerować większej dokładności, niż wynika to z danych.

\section{\ploren{Porównanie wariantów i rozwiązania odniesienia}{Comparison of variants and baseline}}

Ocena rozwiązania powinna obejmować porównanie z wariantem odniesienia: rozwiązaniem istniejącym, metodą bazową, wynikiem analitycznym, danymi producenta albo stanem sprzed wprowadzenia modyfikacji. Warunki porównania muszą być jednakowe albo różnice należy uwzględnić podczas interpretacji.

Nie wystarczy podać zmiany bezwzględnej lub procentowej. Trzeba ocenić, czy różnica ma znaczenie praktyczne, mieści się w granicach niepewności oraz czy została uzyskana kosztem pogorszenia innych właściwości. Dla algorytmu może to być kompromis między dokładnością i czasem obliczeń, a dla urządzenia — między sprawnością, kosztem, masą i niezawodnością.

\section{\ploren{Ocena wymagań i celów}{Evaluation against requirements and objectives}}

Należy wrócić do wymagań i kryteriów zdefiniowanych w rozdziale dotyczącym celu oraz projektu. Dla każdego kluczowego wymagania trzeba wskazać dowód jego spełnienia: wynik testu, pomiar, analizę albo obserwację. Pomocna może być zwięzła tabela zawierająca identyfikator wymagania, kryterium akceptacji, uzyskany wynik i ocenę spełnienia.

Ocena nie powinna być zawyżana. Wymaganie należy uznać za częściowo spełnione lub niespełnione, jeżeli wyniki tego dowodzą, a następnie wyjaśnić przyczynę i konsekwencje. Również wynik negatywny jest wartościowy, jeśli został uzyskany poprawną metodą i prowadzi do istotnego wniosku.

\section{\ploren{Analiza statystyczna i niepewność}{Statistical analysis and uncertainty}}

Jeżeli charakter danych tego wymaga, należy podać miary rozrzutu, przedziały ufności, wyniki testów statystycznych lub analizę niepewności. Wartość \emph{p} nie zastępuje informacji o wielkości efektu i jego znaczeniu technicznym. Zastosowana metoda statystyczna musi odpowiadać rodzajowi danych, liczbie obserwacji i przyjętym założeniom.

W pracach pomiarowych wynik należy przedstawiać wraz z właściwą jednostką i niepewnością. Na wykresach można zastosować słupki błędów albo pasma ufności, podając w podpisie, co dokładnie oznaczają. Jeżeli niepewność nie została oszacowana, nie należy formułować wniosków opartych na różnicach porównywalnych z dokładnością aparatury lub rozrzutem pomiarów.

\section{\ploren{Interpretacja wyników}{Interpretation of results}}

Dyskusja powinna wyjaśniać mechanizmy i przyczyny obserwowanych rezultatów. Należy wskazać, które wyniki były oczekiwane, które zaskakujące oraz czy potwierdzają przyjęte hipotezy lub założenia. Interpretacja musi wynikać z danych; nie należy przedstawiać przypuszczenia jako udowodnionego faktu.

Warto rozważyć alternatywne wyjaśnienia, wpływ zakłóceń i zależności pomiędzy parametrami. Jeżeli zachowanie rozwiązania różni się w poszczególnych warunkach, należy opisać granice jego poprawnego działania zamiast ograniczać się do wyniku średniego.

\section{\ploren{Odniesienie do literatury}{Comparison with the literature}}

Uzyskane wyniki trzeba zestawić z publikacjami omówionymi w przeglądzie literatury. Zgodność może wzmacniać wiarygodność rezultatów, natomiast rozbieżność wymaga analizy różnic w danych, modelach, aparaturze, parametrach i warunkach eksperymentu. Nie należy porównywać samych liczb, jeśli pochodzą z nieporównywalnych procedur.

Dyskusja powinna wskazywać, co praca wnosi względem istniejących rozwiązań: potwierdzenie metody w nowych warunkach, usprawnienie implementacji, ocenę ograniczeń, nowe dane, konstrukcję urządzenia albo inne praktyczne lub poznawcze osiągnięcie. Twierdzenia o nowości muszą być proporcjonalne do zakresu przeprowadzonego przeglądu.

\section{\ploren{Wyniki nieoczekiwane i negatywne}{Unexpected and negative results}}

Nie należy pomijać poprawnie uzyskanych rezultatów tylko dlatego, że nie potwierdzają założeń lub są mniej korzystne od oczekiwanych. Trzeba przedstawić je wraz z analizą możliwych przyczyn. Jeśli wykryto błąd metodyczny lub implementacyjny, należy wskazać, które wyniki pozostają wiarygodne, a których nie można interpretować.

Niepowodzenie określonego wariantu może ujawnić ograniczenie metody i prowadzić do ważniejszego wniosku niż niewielka poprawa wybranego wskaźnika. Nie wolno jednak zastępować brakujących badań spekulacją — kwestie wymagające dalszej weryfikacji należy nazwać wprost.

\section{\ploren{Ograniczenia i zagrożenia dla wiarygodności}{Limitations and threats to validity}}

Należy wskazać czynniki ograniczające pewność wniosków i możliwość ich uogólnienia. Mogą one wynikać z liczby prób, zakresu danych, dokładności aparatury, uproszczeń modelu, doboru parametrów, warunków laboratoryjnych, wykorzystanej platformy albo wpływu autora na ocenę rezultatów.

Ograniczenia powinny być konkretne i powiązane z wynikami. Samo stwierdzenie, że badania przeprowadzono w ograniczonym zakresie, jest niewystarczające. Trzeba wyjaśnić, jakie wnioski pozostają uzasadnione, których nie można rozszerzyć na inne warunki oraz jakie dodatkowe badanie mogłoby zmniejszyć niepewność.

\section{\ploren{Podsumowanie wyników i dyskusji}{Results and discussion summary}}

Na końcu należy zebrać najważniejsze ustalenia w kolejności odpowiadającej celom pracy. Trzeba wskazać stopień spełnienia wymagań, odpowiedzi na pytania badawcze lub wynik weryfikacji hipotez oraz najważniejsze ograniczenia. Podsumowanie tego rozdziału powinno dostarczyć bezpośredniej podstawy do sformułowania końcowych wniosków, ale nie powinno jeszcze powtarzać całego podsumowania pracy.
